% !TeX spellcheck = es_ES
\documentclass[letterpaper,11pt]{article}
\usepackage[utf8]{inputenc}
\usepackage[T1]{fontenc}
\usepackage[spanish]{babel}
\decimalpoint % separador decimal con punto, como en las tablas de la fuente original
\usepackage{lmodern,microtype}
\usepackage{amsmath,amssymb,amsfonts}
\usepackage{geometry}
\geometry{margin=2.54cm}
\usepackage{setspace}
\setstretch{1.15}
\usepackage{enumitem}
\usepackage{booktabs}
\usepackage{color}
\usepackage{hyperref}
\definecolor{ntr}{rgb}{0.8,0,0}
\hypersetup{colorlinks=true,linkcolor=ntr,citecolor=ntr,urlcolor=black}

% Portada (Luis Cabral): titulo \LARGE sans-serif alineado a la izquierda, sin numero de pagina
\makeatletter
    \renewcommand{\@maketitle}{%
    \newpage\parindent0in\null\vskip2em%
    {\LARGE\textbf{\textsf{\@title}}\par\vskip2em}%
    \@author\par\vskip2em%
    \@date%
    \par
    \thispagestyle{empty}\setcounter{page}{1}%
    }%
\makeatother

\DeclareMathOperator{\Var}{Var}
\DeclareMathOperator{\Cov}{Cov}
\DeclareMathOperator{\Corr}{Corr}
\DeclareMathOperator{\E}{\mathbb{E}}
\DeclareMathOperator{\N}{Normal}
\DeclareMathOperator{\LP}{\mathbb{L}}
\newcommand\independent{\protect\mathpalette{\protect\independenT}{\perp}}
\def\independenT#1#2{\mathrel{\rlap{$#1#2$}\mkern2mu{#1#2}}}


\title{Análisis Econométrico \\ Ayudantía 2}
\author{Magíster en Economía, Universidad Alberto Hurtado}
\date{}

\begin{document}
\maketitle

\noindent Esta ayudantía comienza con un control breve (\textbf{Control 1}, 15 minutos, sin
apuntes) sobre el contenido de la Tarea 1. Luego revisamos en conjunto los ejercicios 2, 4, 5 y
6 de la Tarea 1 --- se reproducen aquí para tenerlos a mano durante la revisión.

\vspace{1em}

\noindent\textbf{Ejercicio 2 de la Tarea 1 (10 puntos). Modelos saturados.}

Estamos interesados en el efecto de una variable $x$ en $y$. Suponga que $x$ es una variable
que toma los valores 0, 1 y 2. Defina las siguientes variables aleatorias
\begin{align*}
d_1 =
\begin{cases}
1 & \mbox{si } x = 1, \\
0 & \mbox{en caso contrario, }
\end{cases}
d_2 =
\begin{cases}
1 & \mbox{si } x = 2, \\
0 & \mbox{en caso contrario. }
\end{cases}
\end{align*}
Demuestre que la esperanza condicional $\E(y|d_1,d_2)$ es lineal. \\
Dado lo visto en clase usted puede deducir que la proyección lineal es igual a la esperanza
condicional. \\
El modelo anterior pertenece a una familia de modelos con la característica que la esperanza
condicional es lineal ¿Cuál es esta familia de modelos?

\vspace{1em}

\noindent\textbf{Ejercicio 4 de la Tarea 1 (10 puntos). Agregar variables y el $R^2$.}

Suponga dos modelos de regresión
\begin{align}
y &= x'\gamma+u,\\
y &= x'\beta_1+z'\beta_2+v,
\end{align}
donde el modelo (2) incluye $x$ (las variables explicativas del modelo 1) e incluye $z$
(variables adicionales no incluidas en el modelo 1). Además suponga que $x$ incluye una
constante. Demuestre que $\Var(u)\geq \Var(v)$.\\
Utilice el resultado anterior para argumentar que siempre que agreguemos variables el $R^2$ del
modelo no disminuye.

\vspace{1em}

\noindent\textbf{Ejercicio 5 de la Tarea 1 (20 puntos). Variable omitida.}

Dados los modelos de regresión
\begin{align}
    y &= \gamma_0+\gamma_1 x_{1}+u,\\
    y &= \beta_0 + \beta_1 x_{1}+\beta_2 x_{2}+v.
\end{align}
\begin{enumerate}
    \item Encuentre la relación entre $\beta_1$ y $\gamma_1$. ¿Cuándo se cumple que
    $\beta_1=\gamma_1$?
    \item Suponga que $y = \log(\text{salarios})$, $x_{1}=$ años de educación, y $x_{2}=$
    habilidad innata. Si el modelo (4) mide el efecto causal de educación en salarios, comente
    sobre las consecuencias de utilizar el modelo (3) en lugar del modelo (4).
\end{enumerate}
Pista: Escriba $\gamma_1$ utilizando la fórmula para el modelo de regresión simple, y luego
reemplace $y$ según el modelo (4).

\vspace{1em}

\noindent\textbf{Ejercicio 6 de la Tarea 1 (20 puntos). Error de medida en la variable
dependiente.}

Estamos interesados en el modelo de regresión
\begin{align*}
    q^* &= \beta_0 + \beta_1 x_{1}+r^*
\end{align*}
donde $\E(r^*)=0$ y $\E(x_1\,r^*)=0$. Suponga que no se puede observar $q^*$ pero observamos $q$
tal que $q=q^*+e$ donde $\E(e)=0$ y $\E(x_1\,e)=0$. Se puede interpretar que $q$ es una medición
con errores de medida de $q^*$.
\begin{enumerate}
    \item Dado el modelo lineal de regresión de $q$ (la variable observada con errores de
    medida) en $x_1$:
\begin{align*}
    q &= \delta_0 + \delta_1 x_{1}+error.
\end{align*}
    Demuestre que $\delta_1=\beta_1$ y $\delta_0=\beta_0$.
    \item Suponga que $\E(r^* e)=0$ y demuestre que la varianza del error es mayor en el modelo
    con error de medida.\\
\end{enumerate}

\end{document}
